\cleardoublepage
\section{Rates}

A \textbf{rate} is a special type of ratio that compares two different kinds of units. For example, it can compare miles to hours, dollars to items, or words to minutes.

**Example:**  
If you drive 100 kilometers in 2 hours, the rate is:
\begin{itemize}
  \item 100 km per 2 hours → \textbf{50 km/h} (kilometers per hour)
\end{itemize}

\subsection*{Common Rates}

\begin{itemize}
  \item \textbf{Speed:} miles per hour (mph), kilometers per hour (km/h)
  \item \textbf{Prices:} dollars per item (\$/item)
  \item \textbf{Typing:} words per minute (wpm)
\end{itemize}

\subsection*{Everyday Examples}

\begin{itemize}
  \item A printer prints 60 pages in 3 minutes → Rate: \textbf{20 pages per minute}
  \item A worker earns \$200 for 8 hours → Rate: \textbf{\$25/hour}
\end{itemize}

Rates are everywhere—in recipes, gas mileage (MPG), data usage on your phone, and how fast athletes run.

\newpage
\subsection{Practice Questions}

\begin{enumerate}
  \item \textbf{Short Question} \\
  If a car travels 150 miles in 3 hours, what is the rate in miles per hour?
  
  \textbf{Answer:} \textbf{50 mph} — Divide 150 miles by 3 hours: $150 \div 3 = 50$ mph.

  \item \textbf{Long Question} \\
  A printer can print 240 pages in 8 minutes. What is the printing rate in pages per minute?
  
  \textbf{Answer:} \\
  Rate = $\frac{240 \text{ pages}}{8 \text{ minutes}} = 30$ pages per minute \\
  \textbf{Answer:} 30 pages per minute

  \item \textbf{Challenge Question} \\
  Sarah earns \$180 for working 6 hours, and Tom earns \$120 for working 4 hours. Who has the higher hourly rate?
  
  \textbf{Answer:} \\
  Sarah's rate: $\frac{\$180}{6 \text{ hours}} = \$30$/hour \\
  Tom's rate: $\frac{\$120}{4 \text{ hours}} = \$30$/hour \\
  \textbf{Answer:} They have the same hourly rate of \$30/hour
\end{enumerate}

\newpage
\subsection{Hand-drawing 1}
This is where you will see a hand-drawn answer to the previous question in \textbf{English}.
\begin{figure}[htbp]
  \centering
  \includegraphics[width=\linewidth]{example-image} % TODO: Update caption for actual content
  \caption{A hand-drawn answer in English.}
\end{figure}

\newpage
\subsection{Hand-drawing 2}
This is where you will see a hand-drawn answer to the previous question in \textbf{Korean}.
\begin{figure}[htbp]
  \centering
  \includegraphics[width=\linewidth]{example-image} % TODO: Update caption for actual content
  \caption{A hand-drawn answer in Korean.}
\end{figure}

\cleardoublepage
\subsection*{Short Question (English)}
A car travels 180 kilometers in 3 hours. What is its speed in km/h?

\fullpageimage{images/q3_2_1-eng.jpg}

\newpage
\subsection*{Short Question (Korean)}
A car travels 180 kilometers in 3 hours. What is its speed in km/h?

\fullpageimage{images/q3_2_1-kor.jpg}

\newpage
\subsection*{Short Question (Answer)}
A car travels 180 kilometers in 3 hours. What is its speed in km/h?

\textbf{Answer:} $180 \div 3 = 60$ km/h

\cleardoublepage
\subsection*{Long Question (English)}
A market sells 5 kg of rice for \$20. What is the rate in dollars per kilogram?

\fullpageimage{images/q3_2_2-eng.jpg}

\newpage
\subsection*{Long Question (Korean)}
A market sells 5 kg of rice for \$20. What is the rate in dollars per kilogram?

\fullpageimage{images/q3_2_2-kor.jpg}

\newpage
\subsection*{Long Question (Answer)}
A market sells 5 kg of rice for \$20. What is the rate in dollars per kilogram?

\textbf{Answer:} $20 \div 5 = 4$ → \textbf{\$4 per kg}

\cleardoublepage
\subsection{Challenge Question (English)}
A Korean drama episode is 75 minutes long, and you watch 3 episodes in a day. How many minutes per day do you watch? Express as a rate in hours per day.

\fullpageimage{images/q3_2_3-eng.jpg}

\newpage
\subsection{Challenge Question (Korean)}
A Korean drama episode is 75 minutes long, and you watch 3 episodes in a day. How many minutes per day do you watch? Express as a rate in hours per day.

\fullpageimage{images/q3_2_3-kor.jpg}

\newpage
\subsection{Challenge Question (Answer)}
A Korean drama episode is 75 minutes long, and you watch 3 episodes in a day. How many minutes per day do you watch? Express as a rate in hours per day.

\textbf{Answer:}\\
Total minutes: $75 \times 3 = 225$ minutes\\
Convert to hours: $225 \div 60 = 3.75$ hours/day

